\section*{Bab \ref{ch:b1:reals} --- Bilangan Real}

\begin{solution}{exo:b1:reals:1}
$A$: setiap unsurnya $\leq 1$ dan $1 \in A$: jadi $\sup A = \max A = 1$.
Batas bawahnya: $0$ membatasi dari bawah; dan untuk $\varepsilon > 0$, Archimedes
memasok $n$ dengan $\frac 1n < \varepsilon$, sehingga tak ada bilangan positif yang
membatasi $A$ dari bawah: jadi $\inf A = 0$, yang tak tercapai (tanpa min).

$B$: sukunya $0, -\frac12, \frac23, -\frac34, \frac45, \dots$ Adapun suku
berindeks genap $\frac{n}{n+1}$ (dengan $n$ genap) naik menuju $1$ tanpa mencapainya;
sedangkan suku berindeks ganjil $-\frac{n}{n+1}$ turun menuju $-1$. Jadi $\sup B = 1$ dan
$\inf B = -1$, keduanya tak tercapai: tanpa maks, tanpa min. (Batasnya: $\abs{b}
< 1$ untuk setiap $b \in B$; dan $\frac{n}{n+1} = 1 - \frac{1}{n+1} > 1 -
\varepsilon$ untuk $n$ besar, serupa itu di bawahnya.)

$C = \intoo{-\sqrt 3}{\sqrt 3}$: di sini $\sup C = \sqrt 3$, $\inf C =
-\sqrt 3$, keduanya tak tercapai.
\end{solution}

\begin{solution}{exo:b1:reals:2}
Tulis $x = \lfloor x \rfloor + u$, $y = \lfloor y \rfloor + v$ dengan
$u, v \in \intco{0}{1}$. Maka $x + y = \lfloor x \rfloor + \lfloor y
\rfloor + (u + v)$ dengan $u + v \in \intco{0}{2}$. Jika $u + v < 1$,
$\lfloor x + y\rfloor = \lfloor x\rfloor + \lfloor y \rfloor$; sedangkan jika
$1 \leq u + v < 2$, $\lfloor x+y \rfloor = \lfloor x \rfloor +
\lfloor y \rfloor + 1$. Kedua kasusnya terjadi: karena $(x, y) = (0.2,\, 0.3)$
memberikan kesamaan yang kiri, dan $(0.7,\, 0.8)$ yang kanan.
\end{solution}

\begin{solution}{exo:b1:reals:3}
Misalkan $k = \lfloor x \rfloor$, sehingga $k \leq x < k + 1$. Dengan mengalikannya dengan
$n$: $nk \leq nx < nk + n$, lalu mengambil \omterm{thm:b1:reals:floor}{lantainya} (yaitu operasi yang naik
pada sisi bilangan bulatnya): $nk \leq \lfloor nx \rfloor \leq nk + n
- 1$. Dengan \omterm{def:b1:arith:divides}{membaginya} dengan $n$: $k \leq \frac{\lfloor nx \rfloor}{n} < k + 1$,
sehingga \omterm{thm:b1:reals:floor}{lantai} yang di luarnya adalah $k$.
\end{solution}

\begin{solution}{exo:b1:reals:4}
Setiap unsur $A$ termasuk $B$, sehingga $\sup B$ membatasi $A$ dari atas:
jadi $\sup A \leq \sup B$ (karena $\sup A$ adalah \omterm{def:b1:reals:bounds}{batas atas} yang
\emph{terkecil}). Secara simetris $\inf B \leq \inf A$. Akhirnya $\inf A \leq
\sup A$ karena $A$ tak kosong: yaitu sebarang $a \in A$ duduk di antara keduanya.
\end{solution}

\begin{solution}{exo:b1:reals:5}
Misalkan $s = \sup A$, $t = \sup B$. Setiap $a + b \leq s + t$: jadi itu \omterm{def:b1:reals:bounds}{batas
atas}. Untuk $\varepsilon > 0$, pilihlah $a > s - \frac\varepsilon2$ dan
$b > t - \frac\varepsilon2$ (\cref{prop:b1:reals:epsilon}): maka $a +
b > s + t - \varepsilon$. Menurut pencirian lewat $\varepsilon$,
$\sup(A+B) = s + t$.

Untuk $-A$: di sini $m$ membatasi $-A$ dari atas $\iff$ $-m$ membatasi $A$ dari bawah;
sehingga \omterm{def:b1:reals:bounds}{batas atas} terkecil $-A$ bersesuaian dengan batas bawah
terbesar $A$: jadi $\sup(-A) = -\inf A$.
\end{solution}

\begin{solution}{exo:b1:reals:6}
Untuk setiap $x$: $f(x) + g(x) \leq \sup f + \sup g$; lalu mengambil sup
ruas kirinya memberikan ketaksamaannya. Contoh yang tegas: $E = \{0, 1\}$,
$f = \mathbf{1}_{\{0\}}$ (yang bernilai $1$ di $0$, dan $0$ selainnya), $g =
\mathbf{1}_{\{1\}}$: maka $\sup(f + g) = 1 < 2 = \sup f + \sup g$.

Tak ada pertentangan dengan \cref{exo:b1:reals:5}: karena di sana $a \in A$ dan $b
\in B$ berubah secara \emph{saling bebas}; sedangkan di sini $x$ yang sama menyuapi $f$
maupun $g$ --- jadi \omterm{def:b1:logic:sets}{himpunan} $\{f(x) + g(x) : x \in E\}$ lebih kecil daripada
\omterm{def:b1:logic:sets}{himpunan} $\{f(x) + g(y) : x, y \in E\}$.
\end{solution}

\begin{solution}{exo:b1:reals:7}
Bilangan $\sqrt 6$ irasional: karena $6 = 2 \times 3$ bukan kuadrat sempurna,
dan $v_2(6q^2) = 1 + 2v_2(q)$ yang ganjil mencegah $6q^2 = r^2$ (seperti pada
\cref{exo:b1:arith:7}). Sekarang andaikan $x = \sqrt 2 + \sqrt 3 \in \Q$.
Maka $x^2 = 5 + 2\sqrt 6 \in \Q$, sehingga $\sqrt 6 = \frac{x^2 - 5}{2} \in
\Q$: yang bertentangan. Jadi $\sqrt 2 + \sqrt 3 \notin \Q$.
\end{solution}

\begin{solution}{exo:b1:reals:8}
Misalkan $x < y$. Pilihlah $n \in \N$ dengan $2^n > \frac{1}{y - x}$ (menurut Archimedes:
$2^n \geq n + 1$ lewat induksi yang mudah, sehingga suatu pangkat $2$ melampaui
sebarang bilangan real). Maka, seperti pada bukti \cref{thm:b1:reals:density} dengan
$2^n$ menggantikan $n$: bilangan $m = \lfloor 2^n x \rfloor + 1$ memenuhi $x <
\frac{m}{2^n} < y$. Jadi $D$ padat.
\end{solution}

\begin{solution}{exo:b1:reals:9}
\begin{enumerate}
  \item Andaikan $\alpha > 0$. Pertama, $\alpha \in G$. Andaikan tidak: maka
        menurut pencirian \omterm{def:b1:reals:bounds}{infimum} lewat $\varepsilon$ dengan
        $\varepsilon = \alpha$, ada $g \in G$ dengan $\alpha < g <
        2\alpha$ (tegas di kirinya karena $\alpha \notin G$); lalu,
        dengan $\varepsilon = g - \alpha$, ada $h \in G$ dengan
        $\alpha < h < g$. Sekarang $g - h \in G$ dan $0 < g - h < g -
        \alpha < \alpha$: yaitu sebuah unsur $G \cap \intoo{0}{+\infty}$
        di bawah \omterm{def:b1:reals:bounds}{infimumnya}, yang mustahil. Jadi $\alpha \in G$, dan $\alpha\Z
        \subseteq G$ (karena $G$ sebuah \omterm{def:b1:structures:group}{grup}). Sebaliknya, untuk $x \in G$, misalkan
        $k = \lfloor x/\alpha \rfloor$: maka $x - k\alpha \in G$ dan
        $0 \leq x - k\alpha < \alpha$, dan definisi $\alpha$
        memaksa $x - k\alpha = 0$. Jadi $G = \alpha\Z$.
  \item Andaikan $\alpha = 0$, dan misalkan $x < y$. Ada $g \in G$
        dengan $0 < g < y - x$. Kelipatan $kg$ dengan $k = \lfloor x/g
        \rfloor + 1$ memenuhi $x < kg \leq x + g < y$, dan $kg \in
        G$: jadi padat.
\end{enumerate}
\omterm{def:b1:logic:sets}{Himpunan} $G = \Z + \sqrt 2\,\Z$ merupakan \omterm{def:b1:structures:subgroup}{subgrup} $(\R, +)$. Ia tak berbentuk
$\alpha\Z$: karena kalau tidak, $1 = p\alpha$ dan $\sqrt 2 = q\alpha$
(dengan $p, q \in \Z$) akan memberikan $\sqrt 2 = \frac qp \in \Q$,
yang bertentangan. Menurut dikotominya, $G$ padat di $\R$.
\end{solution}

\begin{solution}{exo:b1:reals:10}
Misalkan $s = \sup A > 0$, $t = \sup B > 0$. Untuk $a \in A$, $b \in B$: $ab
\leq st$ (karena mengalikan ketaksamaan antara bilangan yang \emph{positif}).
Untuk $0 < \varepsilon < \min(s, t)$: pilihlah $a > s - \varepsilon$ dan
$b > t - \varepsilon$; maka
\[
ab > (s - \varepsilon)(t - \varepsilon)
= st - \varepsilon(s + t) + \varepsilon^2
> st - \varepsilon (s + t),
\]
dan $\varepsilon(s+t)$ dapat dibuat sekecil-kecilnya: sehingga menurut pencirian
lewat $\varepsilon$ (dalam bentuk: tak ada bilangan $< st$ yang membatasi
$AB$ dari atas), $\sup AB = st$.

Kepositifannya penting: karena $A = B = \{-1, 0\}$ memberikan $AB = \{0, 1\}$,
$\sup AB = 1$, padahal $\sup A \cdot \sup B = 0 \times 0 = 0$.
\end{solution}

\begin{solution}{exo:b1:reals:11}
Tulis $s = \sup A$, $i = \inf A$. Untuk $a, a' \in A$: dari $a \leq s$
dan $a' \geq i$ diperoleh $a - a' \leq s - i$; lalu menurut kesimetrian $\abs{a -
a'} \leq s - i$, sehingga $s - i$ membatasi \omterm{def:b1:logic:sets}{himpunan} selisihnya dari atas. Untuk
$\varepsilon > 0$, pilihlah $a > s - \frac\varepsilon2$ dan $a' < i
+ \frac\varepsilon2$ (\cref{prop:b1:reals:epsilon} beserta cerminnya
bagi \omterm{def:b1:reals:bounds}{infimum}): maka $\abs{a - a'} \geq a - a' > s - i -
\varepsilon$. Menurut pencirian lewat $\varepsilon$,
$\operatorname{diam} A = s - i$.

Setiap $a \in A$ memenuhi $i \leq a \leq s$, sehingga $A \subseteq
\intcc{i}{s}$, yaitu \omterm{prop:b1:reals:intervals}{selang} tertutup yang panjangnya $\operatorname{diam}
A$. Jika sebuah \omterm{prop:b1:reals:intervals}{selang} tertutup $\intcc{u}{v}$ memuat $A$, maka $v$ menjadi
\omterm{def:b1:reals:bounds}{batas atas} dan $u$ batas bawah $A$, sehingga $u \leq i$ dan $v
\geq s$: jadi $\intcc{i}{s} \subseteq \intcc{u}{v}$. Dengan demikian
$\intcc{i}{s}$ yang terkecil.
\end{solution}

\begin{solution}{exo:b1:reals:12}
Karena $0 \in E$, maka $E \neq \emptyset$. Jika $x > \max(1, y)$ maka $x^2 >
x > y$, sehingga $E$ terbatas di atas oleh $\max(1, y)$: jadi $s = \sup E$
ada (\cref{thm:b1:reals:sup}), dan $s \geq \min(1, y) > 0$
karena $\min(1, y) \in E$: memang jika $y \geq 1$ maka $1^2 = 1
\leq y$, sedangkan jika $y < 1$ maka $y^2 < y$.

\emph{Mustahil bahwa $s^2 < y$.} Pilihlah $0 < h < 1$ dengan $h <
\frac{y - s^2}{2s + 1}$. Maka
\[
(s + h)^2 = s^2 + 2sh + h^2 \leq s^2 + (2s + 1)h < y ,
\]
sehingga $s + h \in E$, yang bertentangan dengan $s$ membatasi $E$ dari atas.

\emph{Mustahil bahwa $s^2 > y$.} Pilihlah $0 < h < s$ dengan $h <
\frac{s^2 - y}{2s}$. Maka $(s - h)^2 = s^2 - 2sh + h^2 > s^2 -
2sh > y$; dan setiap $x \in E$ memenuhi $x^2 \leq y < (s - h)^2$,
sehingga $x < s - h$ (karena keduanya $\geq 0$): jadi $s - h$ menjadi \omterm{def:b1:reals:bounds}{batas atas}
$E$ yang lebih kecil daripada $s$, yang bertentangan dengan keterkecilannya.

Jadi $s^2 = y$. Ketunggalannya: jika $0 < s < s'$ maka $s^2 <
s'^2$, sehingga dua akar positif yang berbeda tak mungkin sama-sama dikuadratkan menjadi $y$.
Kemonotonannya: jika $0 < y < y'$, maka $\sqrt y \neq \sqrt{y'}$, dan
$\sqrt y > \sqrt{y'}$ akan memberikan $y > y'$ lewat pengkuadratan: jadi $\sqrt
y < \sqrt{y'}$.
\end{solution}

\begin{solution}{pb:b1:reals:1}
\textbf{1.} Untuk $\frac 18$: $10 = 8 \cdot 1 + 2$, $20 = 8 \cdot 2
+ 4$, $40 = 8 \cdot 5 + 0$; jadi angkanya $1, 2, 5$, dengan sisa $0$, lalu
nol semuanya: sehingga $\frac 18 = 0.125$. Untuk $\frac 17$: $10 = 7 \cdot 1 +
3$, $30 = 7 \cdot 4 + 2$, $20 = 7 \cdot 2 + 6$, $60 = 7 \cdot 8 +
4$, $40 = 7 \cdot 5 + 5$, $50 = 7 \cdot 7 + 1$: jadi angkanya $1, 4, 2,
8, 5, 7$, dengan sisa $3, 2, 6, 4, 5, 1$. Sisanya sudah
kembali ke $r = 1$, sehingga keenam langkahnya berulang persis selamanya:
$\frac 17 = 0.(142857)$, dengan sisa yang berdaur melewati $1, 3, 2, 6,
4, 5$.

\textbf{2.} Untuk $\frac 13$ dalam basis $2$ (dengan $r_0 = 1$): $2 = 3 \cdot 0 +
2$, $4 = 3 \cdot 1 + 1$, lalu $r = 1$ terulang: jadi $\frac 13 =
(0.\overline{01})_2$. Untuk $\frac{5}{16}$ dalam basis $2$: $10 = 16 \cdot
0 + 10$, $20 = 16 \cdot 1 + 4$, $8 = 16 \cdot 0 + 8$, $16 = 16
\cdot 1 + 0$: jadi $\frac{5}{16} = (0.0101)_2$, yang berakhir. Untuk $\frac 12$
dalam basis $3$: $3 = 2 \cdot 1 + 1$, lalu $r = 1$ terulang seketika:
jadi $\frac 12 = (0.\overline{1})_3$.

\textbf{3.} Sisa algoritmanya setelah $N$ langkah adalah $r_N =
b^N p \bmod q$ (yang dibuktikan secara resmi pada pertanyaan 14; di sini ia
pengamatan bahwa setiap langkahnya mengalikan sisanya dengan $b$ lalu
mereduksinya modulo $q$). Semua angka berikutnya bernilai $0$ jika dan hanya jika suatu
$r_N = 0$, yakni $q \mid b^N p$; dan karena $\gcd(p, q) = 1$, lema
Gauss memberikan $q \mid b^N$. Jika $q \mid b^N$, maka setiap faktor prima
$q$ \omterm{def:b1:arith:divides}{membagi} $b^N$, sehingga \omterm{def:b1:arith:divides}{membagi} $b$ (menurut keprimaannya). Sebaliknya, jika
setiap prima $q = p_1^{a_1} \cdots p_r^{a_r}$ \omterm{def:b1:arith:divides}{membagi} $b$, maka
dengan $A = \max_i a_i$ setiap $p_i^{a_i}$ \omterm{def:b1:arith:divides}{membagi} $b^A$, dan
$p_i^{a_i}$ itu koprima berpasangan, sehingga $q \mid b^A$. Untuk $q = 20 =
2^2 \cdot 5$: kedua primanya \omterm{def:b1:arith:divides}{membagi} $10$ (jadi $\frac{1}{20} = 0.05$),
tetapi $2 \nmid 3$, sehingga $\frac{1}{20}$ berulang selamanya dalam basis $3$.

\textbf{4.} Dari $1^2 = 1 < 2 < 4 = 2^2$ diperoleh $s_0 = 1$. Lalu $1.4^2
= 1.96 < 2 < 2.25 = 1.5^2$: sehingga $\lfloor 10\sqrt 2 \rfloor = 14$,
$s_1 = 1.4$. Berikutnya $1.41^2 = 1.9881 < 2 < 2.0164 = 1.42^2$: $s_2 =
1.41$; $1.414^2 = 1.999396 < 2 < 2.002225 = 1.415^2$: $s_3 =
1.414$; $1.4142^2 = 1.99996164 < 2 < 2.00024449 = 1.4143^2$: $s_4
= 1.4142$. Pada setiap kasusnya ketaksamaan yang terpampang mengatakan persis
$s_n \leq \sqrt 2 < s_n + 10^{-n}$, yang merupakan definisi
\omterm{thm:b1:reals:floor}{lantai} $10^n \sqrt 2$.

\textbf{5.} Di sini $A_0 = \lfloor x \rfloor = 0$ karena $0 \leq x < 1$.
Dari $A_{n-1} \leq b^{n-1} x < A_{n-1} + 1$, kalikan dengan $b$:
\[
b\,A_{n-1} \leq b^n x < b\,A_{n-1} + b .
\]
Bilangan bulat $b\,A_{n-1}$ bernilai $\leq b^n x$, sehingga $b\,A_{n-1} \leq
A_n$; sedangkan $b^n x < b\,A_{n-1} + b$ dengan $b\,A_{n-1} + b$ yang bulat
memaksa $A_n \leq b\,A_{n-1} + b - 1$. Jadi $0 \leq d_n =
A_n - b\,A_{n-1} \leq b - 1$: yaitu sebuah angka.

\textbf{6.} Secara teleskopis: $d_k b^{-k} = A_k b^{-k} - A_{k-1}
b^{-(k-1)}$, sehingga
\[
\sum_{k=1}^n d_k b^{-k} = A_n b^{-n} - A_0 = s_n .
\]
Lalu \omterm{def:b1:arith:divides}{membagi} $A_n \leq b^n x < A_n + 1$ dengan $b^n$ memberikan $s_n \leq x <
s_n + b^{-n}$.

\textbf{7.} Lewat induksi: $b^0 = 1 \geq 1$, dan $b^{n+1} = b \cdot
b^n \geq 2(n + 1) \geq n + 2$. Kemonotonannya: $s_n - s_{n-1} = d_n
b^{-n} \geq 0$. Setiap $s_n \leq x$ (menurut pertanyaan 6): jadi $x$ merupakan \omterm{def:b1:reals:bounds}{batas
atas} $\{s_n\}$. Untuk $\varepsilon > 0$, \omterm{thm:b1:reals:archimedes}{sifat Archimedes}
memasok $n$ dengan $n + 1 > \frac1\varepsilon$, sehingga
$b^{-n} < \varepsilon$, dan lalu $s_n > x - b^{-n} > x -
\varepsilon$ menurut pertanyaan 6. Menurut \cref{prop:b1:reals:epsilon}, $x =
\sup_n s_n$.

\textbf{8.} Andaikan $d_k = b - 1$ untuk setiap $k > N$. Untuk $n > N$,
jumlah geometri yang hingga memberikan
\[
s_n = s_N + (b - 1)\sum_{k=N+1}^{n} b^{-k}
= s_N + b^{-N} - b^{-n} .
\]
Jadi $x \geq s_n = s_N + b^{-N} - b^{-n}$ untuk setiap $n$; lalu dengan membiarkan
suku terakhirnya menyusut di bawah sebarang $\varepsilon$ (pertanyaan 7), $x
\geq s_N + b^{-N}$. Padahal pertanyaan 6 pada peringkat $N$ mengatakan $x < s_N +
b^{-N}$: yang bertentangan. Jadi untaian $(d_n)$ bersifat sejati.

\textbf{9.} Keterbatasannya: $t_n \leq (b-1)\sum_{k=1}^n b^{-k} = 1 -
b^{-n} < 1$, dan $(t_n)$ tidak turun, sehingga $y = \sup t_n$
ada dengan $0 \leq y \leq 1$. Tetapkan $n$. Untuk taksiran
dua sisinya: $t_n \leq y$ sudah jelas. Menurut kesejatiannya pilihlah $m > n$ dengan
$e_m \leq b - 2$. Untuk $p \geq m$:
\[
t_p - t_n = \sum_{k=n+1}^{p} e_k b^{-k}
\leq (b^{-n} - b^{-p}) - b^{-m} < b^{-n} - b^{-m},
\]
karena jumlah di tengahnya kehilangan sekurang-kurangnya $b^{-m}$ terhadap maksimum
yang seluruhnya $(b-1)$; sedangkan untuk $p \leq m$, $t_p \leq t_m \leq t_n + b^{-n} -
b^{-m}$ juga (menurut kemonotonannya ditambah kasus $p = m$). Jadi setiap
$t_p \leq t_n + b^{-n} - b^{-m}$, sehingga $y \leq t_n + b^{-n} -
b^{-m} < t_n + b^{-n}$. (Dengan $n = 0$: $y < 1$, sehingga $y \in
\intco{0}{1}$.) Sekarang $b^n t_n = \sum_{k \leq n} e_k b^{n-k}$ merupakan
bilangan bulat, dan $b^n t_n \leq b^n y < b^n t_n + 1$: sehingga $\lfloor b^n
y \rfloor = b^n t_n$. Akhirnya angka $y$: $d_n(y) = b^n
t_n - b \cdot b^{n-1} t_{n-1} = b^n(t_n - t_{n-1}) = e_n$.

\textbf{10.} Pertanyaan 9 mengatakan: bahwa (nilai untaian) berangka (untaian
itu); sedangkan pertanyaan 5--8 mengatakan: bahwa (angka $x$) membentuk untaian sejati
yang pemenggalannya bersupremum $x$ (pertanyaan 7). Jadi kedua \omterm{def:b1:logic:map}{pemetaannya}
tersusun menjadi identitas dalam kedua urutannya: yakni keduanya bijeksi
yang saling invers antara $\intco{0}{1}$ dan untaian yang sejati. Seandainya dua
untaian sejati bernilai sama, menerapkan \omterm{def:b1:logic:map}{pemetaan} angkanya akan membuat
keduanya sama: jadi tunggal. Inilah \emph{teorema ekspansi
$b$-adik} itu.

\textbf{11.} Misalkan $e_n = b - 1$ untuk $n > M$, dengan $M \geq 0$ yang minimal.
Seperti pada pertanyaan 8, $t_n = t_M + b^{-M} - b^{-n}$ untuk $n \geq M$,
sehingga nilainya $\sup t_n = t_M + b^{-M}$. Jika $M = 0$ maka nilainya
$0 + 1 = 1$: jadi dalam basis $10$, $0.999\dots = 1$ persis --- bukan
hampiran. Jika $M \geq 1$, keminimalannya memberikan $e_M \leq b - 2$,
dan nilainya
\[
t_M + b^{-M} = \frac{b^M t_M + 1}{b^M} \in \intoo{0}{1},
\]
yaitu sebuah pecahan $b$-adik, yang ekspansi \emph{sejatinya} adalah $e_1 \dots
e_{M-1}\,(e_M + 1)\,000\dots$ (karena untaian yang berakhir bersifat sejati,
dan nilainya bilangan yang sama). Sebaliknya bilangan real dengan dua
penyajian pastilah mempunyai satu yang tak sejati (karena kesejatian memakukan
penyajiannya, pertanyaan 10), sehingga berbentuk demikian. Dan setiap
$m/b^N \in \intoo{0}{1}$, yang ditulis dengan angka taknol terakhirnya $d_N$,
memang mempunyai kembaran tak sejati $d_1 \dots d_{N-1}(d_N -
1)(b-1)(b-1)\dots$: jadi pecahan $b$-adiklah yang persis membawa dua
nama, sedangkan bilangan real lainnya hanya satu.

\textbf{12.} Katakanlah untaiannya bersesuaian sampai $m - 1$, dengan pemenggalan
bersama $P = s_{m-1}$, dan $d_m < e_m$. Menurut pertanyaan 9 (yaitu taksiran
atas yang tegas pada peringkat $m$), $x < P + d_m b^{-m} + b^{-m} = P +
(d_m + 1)b^{-m} \leq P + e_m b^{-m} \leq y$, dengan langkah terakhirnya
karena $P + e_m b^{-m}$ adalah pemenggalan $y$, yaitu $t_m \leq y$. Jadi
$d_m < e_m \implies x < y$; lalu dengan menukar perannya, $e_m < d_m \implies
y < x$; dan karena untaiannya berbeda pada $m$, salah satu dari keduanya
berlaku. Kedua arahnya menyusul.

\textbf{13.} Misalkan $z = bx - A_1 \in \intco{0}{1}$ (memang $A_1
\leq bx < A_1 + 1$). Untuk $n \geq 0$: $b^n z = b^{n+1} x - b^n
A_1$ dengan $b^n A_1 \in \Z$, sehingga menurut $\lfloor u - K \rfloor =
\lfloor u \rfloor - K$ (dengan $K$ bulat),
\[
A_n(z) = A_{n+1}(x) - b^n A_1(x) .
\]
Jadi $d_n(z) = A_n(z) - b\,A_{n-1}(z) = A_{n+1} - b^n A_1 -
b\,A_n + b^n A_1 = d_{n+1}(x)$. Sehingga bagian pecahan $bx$
membawa angka yang tergeser; lalu dengan mengulanginya $m$ kali, bagian
pecahan $b^m x$ berangka $(d_{n+m})_{n \geq 1}$.

\textbf{14.} Pembagian Euclid: $b^n p = q\,Q_n + r_n$ dengan $0
\leq r_n < q$. Bagilah dengan $q$: $b^n x = Q_n + \frac{r_n}{q}$ dengan
$0 \leq \frac{r_n}{q} < 1$, sehingga $Q_n = \lfloor b^n x \rfloor =
A_n$, yang memberikan $A_n = \frac{b^n p - r_n}{q}$. Untuk rekurensinya:
$b^n p = b(q\,A_{n-1} + r_{n-1}) = q\,(b\,A_{n-1}) +
b\,r_{n-1}$, sehingga $b^n p$ dan $b\,r_{n-1}$ berselisih kelipatan
$q$: jadi $r_n = (b\,r_{n-1}) \bmod q$.

\textbf{15.} Bagilah $b\,r_{n-1}$ dengan $q$: $b\,r_{n-1} = q\,c +
r_n$ dengan $c = \lfloor b\,r_{n-1}/q \rfloor$. Dengan menyulihkannya ke
tampilan pertanyaan 14: $b^n p = q(b\,A_{n-1} + c) + r_n$,
lalu ketunggalan pembagian Euclid mengenali $A_n = b\,A_{n-1}
+ c$, yakni $d_n = c = \lfloor b\,r_{n-1}/q \rfloor$. Jadi angka
ke-$n$ bergantung pada $r_{n-1}$ saja --- yaitu daur pembagian panjang pada Bagian
I, yang kini bersertifikat.

\textbf{16.} Sebanyak $q + 1$ sisa $r_0, \dots, r_q$ bernilai
dalam \omterm{def:b1:logic:sets}{himpunan} beranggota $q$ yaitu $\intint{0}{q-1}$: sehingga menurut prinsip sarang
merpati (\cref{cor:b1:counting:pigeonhole}) ada dua yang berimpit, katakanlah
$r_N = r_{N+T}$ dengan $0 \leq N < N + T \leq q$. Karena $r_n$
menentukan $r_{n+1}$ (pertanyaan 14), induksi memberikan $r_{n+T} =
r_n$ untuk setiap $n \geq N$; dan karena $r_{n-1}$ menentukan $d_n$
(pertanyaan 15), $d_{n+T} = d_n$ untuk setiap $n \geq N + 1$. Jadi ekspansi setiap
bilangan rasional bersifat periodik pada akhirnya, dengan praperiode $\leq
q$ dan periode $\leq q$.

\textbf{17.} Angka bagian pecahan $b^T y$ adalah
$(d_{n+T}) = (d_n)$ (menurut lema geseran, lalu keperiodikan murninya): jadi
untaian sejati yang sama dengan $y$. Menurut pertanyaan 10 nilainya sama:
$b^T y - \lfloor b^T y \rfloor = y$, sehingga $(b^T - 1)\,y = \lfloor
b^T y \rfloor = A_T \in \N$ dan
\[
y = \frac{A_T}{b^T - 1} ,
\]
yaitu rasional dengan penyebut yang \omterm{def:b1:arith:divides}{membagi} $b^T - 1$; adapun pembilangnya $A_T$
merupakan bilangan bulat yang angka basis $b$-nya $d_1 \dots d_T$. Periksa:
$0.(142857) = \frac{142857}{999999}$, dan $142857 \times 7 =
999999$, sehingga ini $\frac 17$.

\textbf{18.} Jika $d_{n+T} = d_n$ untuk $n > N$, maka bagian pecahan
$z$ dari $b^N x$ berangka $(d_{N+n})_{n\geq1}$ (menurut lema geseran),
yang periodik murni; lalu menurut pertanyaan 17, $z \in \Q$. Kemudian $b^N
x = A_N + z$ memberikan $x = (A_N + z)/b^N \in \Q$. Digabung dengan pertanyaan 16:
$x$ rasional $\iff$ ekspansinya periodik pada akhirnya. Ruas kanannya
menyebut basisnya, sedangkan ruas kirinya tidak: jadi keperiodikan dalam satu basis
setara dengan kerasionalan, sehingga setara pula dengan keperiodikan dalam setiap basis.

\textbf{19.} Untuk $x = \frac 1q$, $r_n = b^n \bmod q$. Jika
$\gcd(b, q) = 1$, maka $r_T = r_0 = 1$ jika dan hanya jika $b^T \equiv
1 \pmod q$; dan $T$ semacam itu ada (karena sarang merpati memberikan $b^i \equiv b^j$,
dengan $i < j$, lalu $b$ terbalikkan modulo $q$, sehingga $b^{j-i} \equiv
1$), lalu yang terkecil --- yaitu orde perkaliannya --- membuat
sisanya, sehingga juga angkanya, periodik murni dengan periode $T$.
Tak ada periode yang lebih kecil: karena periode $T'$ akan memberikan $(b^{T'}
- 1)\frac1q \in \N$ (pertanyaan 17), yakni $q \mid b^{T'} - 1$.
Untuk $q = 7$, $b = 10$: $10 \equiv 3$, $10^2 \equiv 2$, $10^3
\equiv 6$, $10^4 \equiv 4$, $10^5 \equiv 5$, $10^6 \equiv 1
\pmod 7$: jadi ordenya $6$, dan memang $\frac 17$ berperiode enam.

\textbf{20.} Untaiannya memuat tak hingga banyak $0$ (karena angka $x^*$
kebanyakan nol), sehingga ia sejati, dan $x^*$ terdefinisi dengan baik
(pertanyaan 9). Andaikan angkanya periodik pada akhirnya dengan periode
$T$ di luar $N$. Tak hingga banyak angkanya sama dengan $1$ (satu per
bilangan segitiga), sehingga suatu $1$ duduk pada posisi $j > N$; lalu
keperiodikannya menaruh $1$ pada setiap posisi $j + kT$: jadi mulai dari $j$
ke depan, jarak antara $1$ yang berurutan paling jauh $T$. Padahal
$1$ itu duduk persis pada bilangan segitiga, yang jarak berurutannya
$\frac{(j+1)(j+2)}{2} - \frac{j(j+1)}{2} = j + 1$ akhirnya melampaui $T$:
yang bertentangan. Jadi ia tak periodik pada akhirnya, sehingga menurut
pertanyaan 18, $x^* \notin \Q$ --- yaitu keirasionalan yang terbaca dari
irama angkanya belaka.

\textbf{21.} Diberikan $x < y$, pertanyaan 7 memasok $n$ dengan $b^{-n}
< y - x$; tetapkan $m = \lfloor b^n x \rfloor + 1$. Maka $b^n x < m
\leq b^n x + 1 < b^n y$, sehingga $x < \frac{m}{b^n} < y$: jadi padat,
untuk setiap basis sekaligus (dan $b = 2$ memulihkan
\cref{exo:b1:reals:8}). Untuk $\frac pq \in \intoo{0}{1}$ dalam basis
$b = q$: angka pertamanya adalah $\lfloor q \cdot \frac pq \rfloor =
p$ dan bagian pecahan $q \cdot \frac pq = p$ adalah $0$: jadi semua
angka berikutnya lenyap, yaitu ekspansi yang berakhir $\frac pq =
(0.p)_q$. Jadi berakhir bergantung pada basisnya; sedangkan keperiodikan ---
yakni kerasionalan --- tidak (pertanyaan 18).

\textbf{22.} Setiap $e_k \in \{1, 2\}$ merupakan angka basis $10$, dan
untaiannya tak pernah berakhir dengan $9$ semuanya: jadi sejati. Nilainya $y$ terletak di
$\intco{0}{1}$ dan berangka persis $(e_k)$ (pertanyaan 9). Tetapkan
$k$: angka ke-$k$ dari $y$ adalah $e_k$, yang dipilih $\neq$ angka ke-$k$
dari $x_k$, sehingga untaian sejati $y$ dan $x_k$ berbeda, jadi
$y \neq x_k$ (pertanyaan 10: karena penyandiannya \omterm{def:b1:logic:inj}{injektif}). Dengan demikian $y$
tak ada di daftarnya: jadi tak ada \omterm{def:b1:logic:map}{pemetaan} $\N^* \to \intco{0}{1}$ yang \omterm{def:b1:logic:inj}{surjektif}.
Bilangan real, tak seperti bilangan rasional, tak dapat didaftar ---
yaitu ketakterbilangan, yang teorinya dikembangkan \cref{ch:b1:topology}.

\textbf{23.} Jika ekspansinya bersesuaian sampai $n$, maka $x$ dan $y$
mempunyai pemenggalan $s_n$ yang sama, dan pertanyaan 6 menaruh keduanya di
$\intco{s_n}{s_n + b^{-n}}$, yaitu \omterm{prop:b1:reals:intervals}{selang} yang panjangnya $b^{-n}$:
sehingga $\abs{x - y} < b^{-n}$. Konversnya: $x = 0.1$ dan $y = 0.0999$
(yang berakhir, sehingga sejati) memenuhi $\abs{x - y} = 10^{-4} <
10^{-3}$, padahal ekspansinya sudah berbeda pada angka yang pertama.
Yang patut disalahkan adalah pecahan $b$-adik pada pertanyaan 11: karena di dekatnya,
gerakan sekecil apa pun membalikkan setiap angka yang terpampang ($0.0999 \to
0.1000$), sebab merekalah persis bilangan real yang di situ
kembaran tak sejatinya mengintai.

\textbf{24.} Dengan $p = 1$, $q = 10$, $b = 2$, $r_0 = 1$: $2 = 10 \cdot
0 + 2$, $4 = 10 \cdot 0 + 4$, $8 = 10 \cdot 0 + 8$, $16 = 10
\cdot 1 + 6$, $12 = 10 \cdot 1 + 2$ --- lalu $r_5 = 2 = r_1$: jadi
sisanya berdaur $(2, 4, 8, 6)$ mulai indeks $1$. Angkanya: $d_1 =
0$, lalu blok berulangnya $d_2 d_3 d_4 d_5 = 0, 0, 1, 1$:
\[
\tfrac{1}{10} = (0.0\overline{0011})_2 ,
\]
dengan praperiode $1$ dan periode $4$. Menurut pertanyaan 3, ekspansi basis $2$
yang berakhir akan menuntut setiap prima $10$ \omterm{def:b1:arith:divides}{membagi} $2$; padahal prima
$5$ menolak. Jadi $0.1$ \emph{tak} tersajikan oleh untaian
biner yang hingga --- sebuah komputer yang menyimpan bit sebanyak berhingga hanya menyimpan
sebuah pemenggalan, dan galat pemenggalan yang menumpuk itulah sebabnya
$0.1 + 0.2$ pada titik mengambang berbeda dari $0.3$ pada bit terakhirnya.

\textbf{25.} (i) Kelengkapan menghasilkan nilainya: $x = \sup s_n$
dan $y = \sup t_n$ (pertanyaan 7 dan 9) --- karena di atas $\Q$ semata,
untaian sejati $\sqrt 2$ tak akan menamai apa pun. (ii) \omterm{thm:b1:reals:archimedes}{Sifat
Archimedes} membuat $b^{-n}$ akhirnya lebih kecil daripada sebarang
$\varepsilon$, yang memaksa pemenggalannya merapat ke
\omterm{def:b1:reals:bounds}{supremumnya} (pertanyaan 7, 21). (iii) Klausa ketunggalan pada
\omterm{thm:b1:reals:floor}{lantainya} mengenali $Q_n = A_n$ pada pertanyaan 14 dan mengesahkan setiap
penarikan angka $\lfloor u - K \rfloor = \lfloor u \rfloor - K$
(pertanyaan 13). (iv) Prinsip sarang merpati, yang diterapkan pada sisa
yang berhingga banyak, merupakan satu-satunya mesin keperiodikan (pertanyaan 16).
Moralnya: bahwa untaian sejati menyandikan $\intco{0}{1}$ dengan setia lalu mengubah
kerasionalan menjadi irama yang kasatmata; tetapi penjumlahan untaian angka
menuntut simpanan yang merambat dari tak hingga jauhnya di kanan, sehingga tak ada
aturan bertahap hingga yang menghitung angka pertama sebuah jumlah pun ---
sedangkan antarmuka \omterm{def:b1:reals:bounds}{supremum} pada \cref{thm:b1:reals:sup} menangani
seluruh analisis dengan satu aksioma. Jadi angka merupakan \emph{gambaran}
$\R$ yang megah; sedangkan \omterm{def:b1:reals:bounds}{supremum} adalah \emph{mesinnya}.

\end{solution}
